\documentclass[aspectratio=169]{beamer}
\usetheme{simple}

\input{preamble/preamble.tex}
\input{preamble/preamble_math.tex}
% \input{preamble/preamble_acronyms.tex}

\title{Denoising Diffusion Models}
\subtitle{STATS305B: Applied  Statistics II}
\author{Scott Linderman}
\date{\today}


\begin{document}


\maketitle


\begin{frame}{Last Time...}
    \textbf{Outline:}
    \begin{itemize}
        \item Recurrent Neural Networks
        \item Backpropagation Through Time
        \item Vanishing Gradients and Gated RNNs
        \item Other Variations and Uses of RNNs
        \item Revisiting HMMs
        \item Linear RNNs and Parallel Inference
    \end{itemize}
    \vspace{1em}
\end{frame}


\begin{frame}{Today...}
    \textbf{Outline:}
    \begin{itemize}
        \item Denoising Diffusion Models
        \item Noising and Generative Processes
        \item Evidence Lower Bound
        \item Continuous Time Limit
    \end{itemize}
    \vspace{1em}
\end{frame}


\begin{frame}[t,allowframebreaks]{Key Ideas}\label{key-ideas}

Diffusion models work by
\begin{enumerate}
\item  Using a fixed, user-defined \textbf{noising process} to convert data into noise. 
\item Learning to invert this process so that starting from noise, we can
generate samples that approximate the data distribution.
\end{enumerate}

We can think of the DDPM as a giant latent variable model, where the
latent variables are noisy versions of the data. 

As with VAEs, given the latent variables, learning the mapping from latents
to observed data is a supervised regression problem.

\end{frame}

\begin{frame}[t,allowframebreaks]{Noising process}\label{noising-process}

Let \(x \equiv x_0 \in \bbR\) be our observed data (assume scalar for now).  

The noising process is a joint distribution over a sequence of latent variables $x_{0:T}$,
\begin{align*}
q(x_{0:T}) = q(x_0) \prod_{t=1}^T q(x_t \mid x_{t-1}).
\end{align*} 
where $q(x_0) = \frac{1}{n} \sum_{i=1}^n \delta_{x_0^{(i)}}(x_0)$ is the empirical measure of the data.

At each step, the latents will become increasingly noisy
versions of the original data, until at time \(T\) the latent variable
\(x_T\) is essentially pure noise. 

The generative model samples pure noise and attempts to invert the noising
process to produce samples that approximate $q(x_0)$.

\end{frame}

\begin{frame}[t,allowframebreaks]{Gaussian noising process}\label{gaussian-noising-process}

For continuous data, the standard noising process is a first-order
Gaussian autoregressive (AR) process, 
\begin{align*}
q(x_t \mid x_{t-1}) &= \mathrm{N}(x_t \mid \lambda_t x_{t-1}, \sigma_{t}^2).
\end{align*} 

The hyperparameters \(\{\lambda_t, \sigma_t^2\}_{t=1}^T\)
and the number of steps \(T\) are fixed (not learned). We restrict
\(\lambda_t < 1\) so that the process contracts

\end{frame}

\begin{frame}[t,allowframebreaks]{Conditional Distributions}
Since the noising process has linear Gaussian dynamics, we can compute
conditional distributions in closed form.

\begin{align*}
    q(x_{t} \mid x_0) 
    &= \int q(x_{t} \mid x_{t-1}) \, q(x_{t-1} \mid x_0) \dif x_{t-1} \\
    &= \int \mathrm{N}(x_{t} \mid \lambda_{t} x_{t-1}, \sigma_{t}^2) \, \mathrm{N}(x_{t-1} \mid \lambda_{t-1|0} x_0, \, \sigma_{t-1|0}^2) \dif x_{t-1} \\
    &= \mathrm{N}(x_{t} \mid  \lambda_{t} \lambda_{t-1|0} x_0, \lambda_{t}^2 \sigma_{t-1|0}^2 + \sigma_{t}^2) \\
    &= \mathrm{N}\left( x_t \mid \lambda_{t|0} x_0, \sigma_{t|0}^2 \right),
\end{align*}
where the parameters are defined recursively,
\begin{align*}
\lambda_{t|0} &= \lambda_t \lambda_{t-1|0} = \prod_{s=1}^t \lambda_s \\
\sigma_{t|0}^2 &= \lambda_t^2 \sigma_{t-1|0}^2 + \sigma_t^2
\end{align*} 
with base case \(\lambda_{1|0}=\lambda_1\) and
\(\sigma_{1|0}^2=\sigma_1^2\).
\end{frame}

\begin{frame}[t,allowframebreaks]{Variance preserving
diffusions}\label{variance-preserving-diffusions}

It is common to set, 
\begin{align*}
\sigma_t^2 &= 1-\lambda_t^2,
\end{align*} 
in which case the conditional variance simplifies to,
\begin{align*}
\sigma_{t|0}^2 = 1 - \prod_{s=1}^t \lambda_s^2 = 1 - \lambda_{t|0}^2.
\end{align*} 
Under this setting, the noising process preserves the
variance of the marginal distributions. 

If \(\bbE[x_0] = 0\) and
\(\Var[x_0] = 1\), then the marginal distribution of \(x_t\) will be
zero mean and unit variance as well.
\end{frame}

\begin{frame}[t,allowframebreaks]{Limiting Distribution}
Consider the following two limits: 

\begin{enumerate}
    \item As \(T \to \infty\), the conditional distribution goes to a standard normal,
\(q(x_T \mid x_0) \to \mathrm{N}(0, 1)\), which makes the marginal
distribution \(q(x_T)\) easy to sample from.

    \item When \(\lambda_t \to 1\), the noising process adds infinitesimal noise
  so that \(x_t \approx x_{t-1}\), which makes the inverse process
  easier to learn.
\end{enumerate}

These two limits are in conflict with one another! If we add
a small amount of noise at each time step, the inverse process is easier
to learn, but we need to take many time steps to converge to a Gaussian
stationary distribution.

\end{frame}

\begin{frame}[t,allowframebreaks]{Generative process}\label{generative-process}

The generative process is a parameteric model that learns to invert the
noise process, 
\begin{align*}
p(x_{0:T}; \theta) &= p(x_T) \prod_{t=T-1}^0 p(x_t \mid x_{t+1}; \theta).
\end{align*} 

The initial distribution \(p(x_T)\) has no parameters
because it is set to the stationary distribution of the noising process,
\(q(x_\infty)\). 

E.g., for the Gaussian noising process above,
\(p(x_T) = \mathrm{N}(0,1)\).

\end{frame}

\begin{frame}[t,allowframebreaks]{Evidence Lower Bound}\label{evidence-lower-bound}

Like the other latent variable models we studied in this course, we will
estimate the parameters by maximizing an \textbf{evidence lower bound
(ELBO)}, 
\begin{align*}
\cL(\theta) 
&= \E_{q(x_0)} \E_{q(x_{1:T} \mid x_0)} \left[ \log p(x_{0:T}; \theta) - \log q(x_{1:T} \mid x_0) \right] \\
&= \E_{q(x_0)} \E_{q(x_{1:T} \mid x_0)} \left[ \log p(x_{0:T}; \theta)  \right] + c,
\end{align*} 
where \(q(x_{1:T} \mid x_0)\) is the conditional
distibution of \(x_{1:T}\) under the noising process. 

Since $q$ is fixed, the objective simplifies to
\textbf{maximizing the expected log likelihood}.

We can simplify further by expanding the log probability of the
generative model, 
\begin{align*}
\cL(\theta)
&= \E_{q(x_0)} \sum_{t=0}^{T-1} \E_{q(x_t, x_{t+1} \mid x_0)} \left[  \log p(x_t \mid x_{t+1}; \theta) \right] \\
&\propto \E_{q(x_0)} \mathbb{E}_{t \sim \mathrm{Unif}(0,T-1)} \E_{q(x_t, x_{t+1} \mid x_0)} \left[ \log p(x_t \mid x_{t+1}; \theta) \right]
\end{align*} 
which only depends on pairwise conditionals.
\end{frame}

\begin{frame}[t,allowframebreaks]{Gaussian generative
process}\label{gaussian-generative-process}

Since the noising process above adds a small amount of Gaussian noise at
each step, it is reasonable to model the generative process as Gaussian
as well, 
\begin{align*}
p(x_t \mid x_{t+1}; \theta) 
&= 
\mathrm{N}(x_t \mid \mu_\theta(x_{t+1}, t), \widetilde{\sigma}_t^2)
\end{align*} 
where 
\begin{itemize}
    \item \(\mu_\theta: \reals \times [0,T] \mapsto \reals\) is
a nonlinear \textbf{mean function} that should \textbf{denoise}
\(x_{t+1}\) to obtain the expected value of \(x_t\)
    \item \(\widetilde{\sigma}_t^2\) is a fixed variance for the generative
process.
\end{itemize}
\end{frame}

\begin{frame}[t,allowframebreaks]{Parameter sharing}
    
    Rather than learn a separate
function for each time point, it is common to parameterize the mean
function as a function of both the state \(x_{t+1}\) and the time \(t\).

E.g., \(\mu_\theta(\cdot, \cdot)\) can be a neural network that
takes in the state and a positional embedding of the time \(t\), like
the sinusoidal embeddings used in transformers. 
\end{frame}

\begin{frame}[t,allowframebreaks]{Generative Process Variance}

    You could try to learn the
generative process variance as a function of \(x_{t+1}\) and \(t\) as
well, but the literature suggests this is difficult to make work in
practice. 

Instead, is common to set the variance to either 
\begin{itemize}
    \item \(\widetilde{\sigma}_t^2 = \sigma_t^2 = 1-\lambda_t^2\), the conditional
variance in the noising process, which tends to \emph{overestimate} the
conditional variance of the true generative process
    \item  \(\widetilde{\sigma}_t^2 = \Var_q[x_t \mid x_0, x_{t+1}]\), the
conditional variance of the noising process \emph{given} the data
\(x_0\) and the next state \(x_{t+1}\). This tends to
\emph{underestimate} the conditional variance of the true generative
process. 
\end{itemize}

\end{frame}

\begin{frame}[t,allowframebreaks]{Rao-Blackwellization}\label{rao-blackwellization}

Under this Gaussian model for the generative process, we can
analytically compute one of the expectations in the ELBO. This is called
\textbf{Rao-Blackwellization}. It reduces the variance of the objective, which is
good for SGD!

Using the chain rule and the Gaussian generative model,
\begin{align*}
\E_{q(x_t, x_{t+1} \mid x_0)} \left[ \log p(x_t \mid x_{t+1}; \theta) \right] 
&= \E_{q(x_{t+1} \mid x_0)} \E_{q(x_t \mid x_{t+1}, x_0)} \left[\log \mathrm{N}(x_t \mid \mu_\theta(x_{t+1}, t), \widetilde{\sigma}_t^2) \right]
\end{align*}

We already computed the conditional distribution
\(q(x_{t+1} \mid x_0) = \mathrm{N}(x_{t+1} \mid \lambda_{t+1|0} x_0, \sigma_{t+1|0}^2)\)
above. It turns out the second term is Gaussian as well!
\end{frame}

\begin{frame}[t,allowframebreaks]{Conditionals of a Gaussian noising process}

Show that 
\begin{align*}
q(x_t \mid x_{t+1}, x_0) 
&= \mathrm{N}(x_t \mid \mu_{t|t+1,0}, \sigma_{t|t+1,0}^2)
\end{align*} 
where 
\begin{align*}
\mu_{t|t+1,0} &= a_t x_0 + b_t x_{t+1}
\end{align*} 
is a \textbf{linear combination} of \(x_0\) and \(x_{t+1}\)
with weights, 
\begin{align*}
a_t &= \frac{\sigma_{t|t+1,0}^2 \, \lambda_{t|0} }{\sigma_{t|0}^2}  \\
b_t &= \frac{\sigma_{t|t+1,0}^2 \, \lambda_{t+1}}{\sigma_{t+1}^2} \\
\sigma_{t|t+1,0}^2  
&= \left(\frac{1}{\sigma_{t|0}^2} + \frac{\lambda_{t+1}^2}{\sigma_{t+1}^2} \right)^{-1} 
\end{align*}
\end{frame}

\begin{frame}[t,allowframebreaks]{Derivation}
By Bayes rule and the Markovian structure of the noising process, 
\begin{align*}
q(x_t \mid x_{t+1}, x_0) 
&\propto q(x_t \mid x_0) \, q(x_{t+1} \mid x_t) \\
&= \mathrm{N}(x_t \mid \lambda_{t|0} x_0, \sigma_{t|0}^2) \, \mathrm{N}(x_{t+1} \mid \lambda_{t+1} x_t, \sigma_{t+1}^2) \\
&= \mathrm{N}(x_t \mid \mu_{t|t+1,0}, \sigma_{t|t+1,0}^2) 
\end{align*} 
where, by completing the square, \begin{align*}
\sigma_{t|t+1,0}^2 &= \left(\frac{1}{\sigma_{t|0}^2} + \frac{\lambda_{t+1}^2}{\sigma_{t+1}^2} \right)^{-1} \\
\mu_{t|t+1,0} &= \sigma_{t|t+1,0}^2 \left(\frac{\lambda_{t|0} x_0}{\sigma_{t|0}^2} + \frac{\lambda_{t+1} x_{t+1}}{\sigma_{t+1}^2} \right).
\end{align*} 
The forms for \(a_t\) and \(b_t\) can now be read off.
\end{frame}

\begin{frame}[t,allowframebreaks]{Gaussian cross-entropy}

Finally, to simplify the objective we need the Gaussian cross-entropy,

Let \(q(x) = \mathrm{N}(x \mid \mu_q, \sigma_q^2)\) and
\(p(x) = \mathrm{N}(x \mid \mu_p, \sigma_p^2)\). 

Show that,
\begin{align*}
\E_{q(x)}[\log p(x)] &= \log \mathrm{N}(\mu_q \mid \mu_p, \sigma_p^2) -\frac{1}{2} \frac{\sigma_q^2}{\sigma_p^2}
\end{align*} 
\end{frame}

\begin{frame}[t,allowframebreaks]{The Simplified ELBO}
Putting it all together, 
\begin{align*}
\cL(\theta) 
&= \E_{q(x_0)} \E_t \E_{q(x_{t+1} \mid x_0)} \E_{q(x_t | x_0, x_{t+1})} \left[ \log p(x_t \mid x_{t+1}; \theta) \right] \\
&= \E_{q(x_0)} \E_t \E_{q(x_{t+1} \mid x_0)} \left[ \log \mathrm{N}(a_t x_0 + b_t x_{t+1} \mid \mu_\theta(x_{t+1}, t), \widetilde{\sigma}_t^2) -\frac{1}{2} \frac{\sigma_{t|t+1,0}^2}{\widetilde{\sigma}_t^2} \right] \\
&= \frac{1}{2} \E_{q(x_0)} \E_t \E_{q(x_{t+1} \mid x_0)} \left[ \frac{1}{\widetilde{\sigma}_t^2} \left( a_t x_0 + b_t x_{t+1} - \mu_\theta(x_{t+1}, t) \right)^2 \right] + c
\end{align*} 
where we have absorbed terms that are independent of
\(\theta\) into the constant \(c\).
\end{frame}

\begin{frame}[t,allowframebreaks]{Denoising mean function}\label{denoising-mean-function}

The loss function above suggests a particular form of the mean function,
\begin{align*}
\mu_\theta(x_{t+1}, t) &= a_t \hat{x}_0(x_{t+1}, t; \theta) + b_t x_{t+1},
\end{align*} 
where the only part that is learned is
\(\hat{x}_0(x_{t+1}, t; \theta)\), a function that attempts to
\textbf{denoise} the current state. 

Since \(x_{t+1}\) is given and
\(a_t\) and \(b_t\) are determined solely by the hyperparameters, we can
use them in the mean function.

Under this parameterization, the loss function reduces to,
\begin{align*}
\cL(\theta) 
&= \frac{1}{2} \E_{q(x_0)} \E_t \E_{q(x_{t+1} \mid x_0)} \left[ \frac{a_t^2}{\widetilde{\sigma}_t^2} \left(x_0 - \hat{x}_0(x_{t+1}, t; \theta) \right)^2 \right] + c
\end{align*}

One nice thing about this formulation is that the mean function is
always outputting \textit{the same thing} --- an estimate of the completely
denoised data, \(\hat{x}_0\), regardless of the time \(t\).

\end{frame}

\begin{frame}[t,allowframebreaks]{Inverting the noising
process}\label{inverting-the-noising-process}

The generative process attempts to invert the noising process, but what
is the actual inverse of the process? 

Since the noising process is a
Markov chain, the reverse of the noising process must be Markovian as
well. 
\begin{align*}
q(x_{0:T}) &= q(x_T) \prod_{t=T-1}^0 q(x_t \mid x_{t+1})
\end{align*} 
for some sequence of transition distributions
\(q(x_t \mid x_{t+1})\). 

We can obtain those transition distributions by
marginalizing and conditioning, 
\begin{align*}
q(x_t \mid x_{t+1}) 
&= \int q(x_0, x_t \mid x_{t+1}) \dif x_0 \\
&= \int q(x_t \mid x_0, x_{t+1}) \, q(x_0 \mid x_{t+1}) \dif x_0.
\end{align*} 
Using Bayes' rule, 
\begin{align*}
q(x_0 \mid x_{t+1}) 
&= \frac{q(x_0) \, q(x_{t+1} \mid x_0)}{\int q(x_0') q(x_{t+1} \mid x_0') \dif x_0'} 
\end{align*} 
Now recall that
\(q(x_0) = \frac{1}{n} \sum_{i=1}^n \delta_{x_0^{(i)}}(x_0)\) is the
empirical measure of the data \(\{x_0^{(i)}\}_{i=1}^n\). Using this
fact, the conditional is, 
\begin{align*}
q(x_0^{(i)} \mid x_{t+1}) 
&= \frac{q(x_{t+1} \mid x_0^{(i)})}{\sum_{j=1}^n q(x_{t+1} \mid x_0^{(j)})} 
\triangleq w_i(x_{t+1}),
\end{align*} 
where we have defined the weights \(w_i(x_{t+1})\) for each
data point \(i=1,\ldots,n\). They are non-negative and sum to one.

Finally, we can give a simpler form for the optimal generative process,
\begin{align*}
q(x_t \mid x_{t+1}) 
&= \sum_{i=1}^n w_i(x_{t+1}) \, q(x_t \mid x_0^{(i)}, x_{t+1}) \\
&= \sum_{i=1}^n w_i(x_{t+1}) \, \mathrm{N}(x_t \mid a_t x_0^{(i)} + b_t x_{t+1}, \sigma_{t|t+1,0}^2),
\end{align*} 
which we recognize as a mixture of Gaussians, all with the
same variance, with means biased toward each of the \(n\) data points,
and weighted by the relative likelihood of \(x_0^{(i)}\) having produced
\(x_{t+1}\).

For small step sizes, that mixture of Gaussians can be approximated by a
single Gaussian with mean equal to the expected value of the mixture,
\begin{align*}
\E[x_t \mid x_{t+1}]
&= \sum_{i=1}^n w_i(x_{t+1}) \, \left(a_t x_0^{(i)} + b_t x_{t+1} \right) 
\end{align*} 
For small steps, this expected value is approximately,
\begin{align*}
\E[x_t \mid x_{t+1}]
&\approx \frac{x_{t+1}}{\lambda_{t+1}} + \sigma_{t+1}^2 \sum_{i=1}^n w_i(x_{t+1}) \left(\frac{ \lambda_{t|0} x_0^{(i)} - x_{t+1}}{\sigma_{t|0}^2} \right)
\end{align*}
(See online notes for derivation.)

% :::\{admonition\} Derivation :class: dropdown Using the definitions from
% above, \begin{align*}
% a_t &= \frac{\sigma_{t|t+1,0}^2 \, \lambda_{t|0} }{\sigma_{t|0}^2}  \\
% b_t &= \frac{\sigma_{t|t+1,0}^2 \, \lambda_{t+1}}{\sigma_{t+1}^2} 
% \end{align*} we can write \begin{align*}
% \E[x_t \mid x_{t+1}]
% &= \frac{\sigma_{t|t+1,0}^2  \lambda_{t+1}}{\sigma_{t+1}^2} x_{t+1} + \sum_{i=1}^n w_i(x_{t+1}) \frac{\sigma_{t|t+1,0}^2 \lambda_{t|0}}{\sigma_{t|0}^2} x_0^{(i)} 
% \end{align*} Now add and subtract
% \(\frac{\sigma_{t|t+1,0}^2 x_{t+1}}{\sigma_{t|0}^2}\) to obtain,
% \begin{align*}
% \E[x_t \mid x_{t+1}]
% &= \left(\frac{\sigma_{t|t+1,0}^2 \lambda_{t+1}}{\sigma_{t+1}^2} + \frac{\sigma_{t|t+1,0}^2}{\sigma_{t|0}^2}\right) x_{t+1} + \sigma_{t|t+1,0}^2 \sum_{i=1}^n w_i(x_{t+1}) \left(\frac{ \lambda_{t|0} x_0^{(i)} - x_{t+1}}{\sigma_{t|0}^2} \right).
% \end{align*} In the limit of many small steps where
% \(\sigma_{t+1} \to 0\) so that \(\sigma^2_{t+1} \ll \sigma^2_{t|0}\), we
% have \begin{align*}
% \frac{\sigma_{t|t+1,0}^2 \lambda_{t+1}}{\sigma_{t+1}^2} + \frac{\sigma_{t|t+1,0}^2}{\sigma_{t|0}^2}
% \approx \frac{1}{\lambda_{t+1}}
% \end{align*} and \begin{align*}
% \sigma_{t|t+1,0}^2 \approx \sigma_{t+1}^2.
% \end{align*} These approximations complete the derivation. :::

Though it's not immediately obvious, the second term in the expectation
is related to the \textbf{marginal probability}, 
\begin{align*}
q(x_t) 
&= \frac{1}{n} \sum_{i=1}^n q(x_t \mid x_0^{(i)}) \\
&= \frac{1}{n} \sum_{i=1}^n \mathrm{N}(x_t \mid \lambda_{t|0} x_0^{(i)}, \sigma_{t|0}^2) 
\end{align*} 

Specifically, the second term is the \textbf{Stein score
function} of the marginal probability, 
\begin{align*}
\nabla_{x} \log q_t(x_{t+1})
&= \frac{\nabla_{x} q_t(x_{t+1})}{q_t(x_{t+1})}  \\
&=\frac{\frac{1}{n} \sum_{i=1}^n \mathrm{N}(x_{t+1} \mid \lambda_{t|0} x_0^{(i)}, \sigma_{t|0}^2) \left(-\frac{(x_{t+1} - \lambda_{t|0}x_0^{(i)})}{\sigma_{t|0}^2} \right)}{\frac{1}{n} \sum_{j=1}^n \mathrm{N}(x_{t+1} \mid \lambda_{t|0} x_0^{(j)}, \sigma_{t|0}^2) } \\
&= \sum_{i=1}^n w_i(x_{t+1}) \left(\frac{\lambda_{t|0}x_0^{(i)} - x_{t+1}}{\sigma_{t|0}^2} \right)
\end{align*}
\end{frame}

\begin{frame}[t,allowframebreaks]{Final Form}
Putting it all together, for
small steps, the reverse process is approximately Gaussian with mean and
variance, 
\begin{align*}
\E[x_t \mid x_{t+1}]
&\approx \frac{x_{t+1}}{\lambda_{t+1}} + \sigma_{t+1}^2 \nabla_x \log q_t(x_{t+1})  \\
\Var[x_t \mid x_{t+1}] &\approx \sigma_{t+1}^2.
\end{align*} 
This has a nice interpretation: to invert the noise
process, \textbf{first undo the contraction and then take a step in the
direction of the Stein score!} 
\end{frame}

\begin{frame}[t,allowframebreaks]{Continuous time limit}\label{continuous-time-limit}

In practice, the best performing diffusion models are based on a
continuous-time formulation of the noising process as an SDE (Song et al., 2020). 

To motivate this approach, think of
the noise process above as a discretization of a continuous process
\(x(t)\) for \(t \in [0,1]\) with time steps of size
\(\Delta = \tfrac{1}{T}\). 

That is, map \(x_i \mapsto x(i/T)\),
\(\lambda_i \mapsto \lambda(i/T)\), and \(\sigma_i \mapsto \sigma(i/T)\)
for \(i=0,1,\ldots, T\). 

Then the discrete model is can be rewritten as,
\begin{align*}
x(t + \Delta) 
&\sim \mathrm{N}(\lambda(t) x(t), \sigma(t)^2),
\end{align*} 
or equivalently, 
\begin{align*}
x(t + \Delta) - x(t) 
&\sim \mathrm{N}\left(f(x(t), t) \Delta, g(t)^2 \Delta \right)  \\
f(x, t) &= \frac{1 - \lambda(t)}{\Delta} x \\
g(t) &= \frac{\sigma(t)}{\Delta}.
\end{align*} 
We can view this as a discretization of the SDE,
\begin{align*}
\dif X &= f(x, t) \dif t + g(t) \dif W
\end{align*} where \(f(x,t)\) is the \textbf{drift} term, \(g(t)\) is
the \textbf{diffusion} term, and \(\dif W\) is the \textbf{Brownian
motion}.

The reverse (generative) process can be cast as an SDE as well!

Following our derivation of the inverse process above, we can show that
the reverse process is, 
\begin{align*}
\dif X = \left[f(x, t) - g(t)^2 \nabla_x \log q_t(x)\right] \dif t + g(t) \dif W
\end{align*} 
where \(\dif t\) is a \textbf{negative} time increment and
\(\dif W\) is Brownian motion run in reverse time.

\end{frame}

\begin{frame}[t,allowframebreaks]{Multidimensional models}\label{multidimensional-models}

Very few things need to change in order to apply this idea to
multidimensional data \(\mbx_0 \in \reals^D\). 

The standard setup is to
apply a Gaussian noising process to each coordinate \(x_{0,d}\)
independently. 

Then, in the generative model, \begin{align*}
p(\mbx_t \mid \mbx_{t+1}; \theta)
&= \prod_{d=1}^D p(x_{t,d} \mid \mbx_{t+1}; \theta) \\
&= \prod_{d=1}^D \mathrm{N}(x_{t,d} \mid \mu_\theta(\mbx_{t+1}, t, d), \widetilde{\sigma}_t^2). 
\end{align*} 
The generative process still produces a factored
distribution, but we need a separate mean function for each coordinate.

Moreover, the mean function needs to consider the entire state
\(\mbx_{t+1}\). The reason is that \(x_{t,d}\) is not conditionally
independent of \(x_{t+1,d'}\) given \(x_{t+1,d}\); the coordinates are
coupled in the inverse process since all of \(\mbx_{t+1}\) provides
information about the \(\mbx_0\) that generated it.
\end{frame}

\begin{frame}[t,allowframebreaks]{Conclusion}\label{conclusion}

There's a lot we didn't cover! 

The Stein score function that appeared in
the inverse of the noising process allows for connections between
denoising score matching \{cite:p\}\texttt{song2019generative} and
denoising diffusion models.

Another important topic is \textbf{conditional generation}. Suppose we
want to take in text and spit out images, like DALL-E 2 or Stable
Diffusion. One way to do so is using a diffusion model, but to steer the
reverse diffusion based on the text prompt. 

Finally, this class was nominally about models for discrete data, but
this lecture has focused on continuous diffusions. There has been recent
work on discrete denoising diffusion models, which we'll have to cover
another time!

\end{frame}
    
    
\end{document}
